Ch1: \chapter{Heat mortality and the neighbourhood}\label{chapter:problem}

\section{Introduction}

Hot days kill, and the deaths are hard to see. The dying happens
indoors, mostly among people already ill, so each death has a nearer
cause to be filed under. The days themselves are multiplying: the
annual average temperature over Canadian land has risen 1.7\,°C since
1948, roughly double the global increase, and the projections hold
that ratio under every emissions scenario\needcite{CCCR ---confirm
1.7 since 1948 and the double-the-global ratio}.

The death record gives a date, a cause, and a place. It does not give a
line that says heat. The toll is therefore estimated rather than counted:
how many people died across the hot days, against how many would have
been expected to die had those days been cooler.

A city is not one place under heat. Its built surfaces hold the day's
warmth unevenly, and a summer night can run several degrees warmer in one
neighbourhood than in another. The people who meet that night differ too.
Age, illness, income, and access to cooling vary from block to block, and
each changes what a given temperature does.

The comparison is run over an area, and the area cannot be small. It needs
enough deaths across enough hot days for the difference to show against
ordinary variation. Estimated over a whole city, the toll is one number,
and the differences between its neighbourhoods are averaged into it rather
than resolved by it.

\section{Heat mortality as it is measured}

That expectation comes from a curve. Over a place's full range of
temperature, deaths fall to a minimum at one point and rise on either side
of it, gradually into the cold and steeply into the heat. The temperature
at that floor is the minimum-mortality temperature. It sits near the warm
end of what a place is accustomed to, so heat begins to kill at a lower
temperature in a cool city than in a warm one.

\begin{figure}[ht]
\centering
\caption{The temperature--mortality association, drawn schematically.}
\label{fig:erf}
\end{figure}

A hot day keeps killing after it ends. The risk it carries decays as the
days pass; the deaths owed to one hot day are spread across the days that
follow it. This delay is the lag, and the effect of a temperature is that
risk summed over the lag window.

The curve differs between cities, and inside a city, between the people
who live there. A characteristic that changes what a given temperature
does is an effect modifier. Age is among the strongest. Estimates are
therefore made separately by age group rather than pooled across a whole
population.

In Canada the scale has been coming down. Ontario has been split into
fourteen health regions, each fitted separately and pooled to a provincial
curve~\citep{bai2017coronary}. Twelve cities have been fitted one at a
time, returning a single heat effect each~\citep{quick2024twelvecities}.
In Montreal, Toronto, and Vancouver the estimate has been taken below the
city, to the census tract~\citep{quick2025smallarea}. That work covers
three cities and five summers, and names smaller areal units and longer
records as what remains.

A neighbourhood carries too few deaths for a curve to be fitted to it
alone. The curve must therefore be borrowed, and the field has split over
where to borrow it from. One camp borrows from adjacency: the coefficients
of the curve are given a spatial prior, so an area with thin counts is
drawn toward the areas around it~\citep{quijalzamorano2024sbdlnm,
rutten2026spatialvarying,rutten2026effectmod,quick2025smallarea}. The
other borrows from resemblance: curves are fitted where the counts hold
up, pooled against measured characteristics of those places, and read back
at an area from its own~\citep{gasparrini2022smallarea,
masselot2025extensions,vazquez2025norway,dimaso2025italy}. The first is
fitted in one stage, the second in two. Sampling the one-stage model by
Markov chain is slow---seventy-three Barcelona neighbourhoods took
several days---and that cost is what held the Canadian small-area work
to three cities. Approximating the fit instead has since removed it:
391 Sicilian municipalities take
minutes~\citep{rutten2026spatialvarying}.

Those area characteristics are also mapped on their own. A vulnerability
index takes a set of small-area variables, reduces them by principal
components, and sums the components into one score per area. Three themes
recur across these indices---deprivation, age, and isolation---with
measures of the built and green environment entering
alongside~\citep{conlon2020hvipca,mallen2019dallas,sera2019urban,
gasparrini2022smallarea,quick2025smallarea}. In its standard construction
no health outcome enters the score, so it ranks neighbourhoods by
composition rather than by risk.

The score has not been shown to predict mortality. Where indices have been
scored against a mortality outcome they have explained little or none of
it, and the result is cited as grounds for treating them as screening
tools rather than as estimates~\citep{conlon2020hvipca,mallen2019dallas}.
Neither test fitted a temperature--mortality association at the scale it
mapped: one scored the index against the share of deaths falling on hot
days, the other against deaths modelled from a single citywide risk curve.

\section{The apparatus}

The expectation comes from the place itself. Days are grouped into
strata---one neighbourhood, one season, one year---and each day's count is
compared only against the other days in its own stratum. Whatever holds
steady across a stratum drops out of the comparison without ever being
measured. The weather does not, and that is what the model is asked about.
The design is the case time series, and the model fitted to it is a
conditional Poisson regression~\citep{gasparrini2021,armstrong2014}.

The curve is fitted within one city at a time, where the deaths are
numerous enough to support it. The fitted curves differ from city to
city, and the differences track what the cities are: their climate,
the age of their people, the conditions they live in. A second
regression takes the fitted curves themselves as its data points and
estimates that relationship---a meta-regression, with cities as the
units and their measured characteristics as the predictors. A
neighbourhood then enters with the same characteristics, measured at
its own scale, and receives the curve its profile implies. Nothing is
fitted to its deaths; the estimate is borrowed from every city at
once. The apparatus exists at national scale: England and Wales,
mapped to 34{,}753 small areas from curves fitted in 348
districts~\citep{gasparrini2022smallarea}.

\section{Aims}

This thesis resolves heat-related mortality risk to the neighbourhood,
for all of Canada's 41 census metropolitan areas at once. The
instrument is the apparatus just described: curves fitted where the
deaths support them, pooled against what the cities measurably are, and
read back at each dissemination area from its own profile. What is new
is the country, the scale the risk is returned at, and the design
choices the framework leaves open.

Four questions organize the work. How much does the risk vary from
one neighbourhood to the next? What is a neighbourhood made of that
moves its risk? For whom is the risk highest? And can an estimate
borrowed from the national sample be trusted at the scale it is
spent? The first three are what the resolution is for. The fourth is
what it costs: every curve on the map is a prediction, and a
prediction earns its map only by being tested.

Chapter~\ref{chapter:method} builds the apparatus: the within-city
model, the pooling across cities, and the downscale that returns a
curve for every dissemination area. Three chapters then read the
surface it produces. Chapter~\ref{chapter:nationalmap} maps the risk
across all 41 metropolitan areas and answers how much it varies;
Chapter~\ref{chapter:vulnerability} builds the vulnerability profile
and answers what a neighbourhood is made of that moves it;
Chapter~\ref{chapter:agefork} splits the surface by age and answers
for whom it is highest. Chapter~\ref{chapter:validation} takes up the
fourth question, testing the predicted surface against simulation
where the truth is planted and against the record where it is not.
Chapter~\ref{chapter:discussion} weighs what the answers establish.